
\section{Prerequisites}
In the context of MLE agents, the main goal is to maximize the evaluation score on a given benchmark under a limited time budget.  
Formally, let $\mathcal{T}$ be an ML task with dataset $\mathcal{D} = \{\mathcal{D}_{\text{dev}}, \mathcal{D}_{\text{test}}\}$, where $\mathcal{D}_{\text{dev}}$ is the development set and $\mathcal{D}_{\text{test}}$ is the test set. The agent can only develop, tune, and iterate its solution on $\mathcal{D}_{\text{dev}}$, while the final performance is measured on the unseen $\mathcal{D}_{\text{test}}$.
This fundamental separation means the agent never has access to the final evaluation metric during its development process. Instead, it must rely on a proxy metric, $M(s; \mathcal{D}_{\text{dev}})$, evaluated on the development set to guide its exploration. 
Let $T$ denote the total allowed wall-clock time for the agent to generate a complete solution (including the actual execution/running time of the solution code).  
Let $s = \text{Agent}(\mathcal{T}, T)$ be a candidate solution generated by the agent (e.g., a Python script for data preprocessing, model training, and evaluation) within this generation-time budget $T$.  
The performance of $s$ on $\mathcal{T}$ is measured by a task-specific metric $M(s; \mathcal{D}_{\text{test}}) \in \mathbb{R}$, such as accuracy or correlation.  


The agent’s objective can be written as the constrained optimization problem  
\[
s^\star = \arg\max_{\substack{s = \text{Agent}(\mathcal{T}, T),\\ s \in \mathcal{S}}} M(s; \mathcal{D}_{\text{test}}), \quad \text{s.t.} \quad \text{gen\_time}(s) \le T,
\]  
where $\mathcal{S}$ is the set of all feasible solutions the agent can generate and $\text{gen\_time}(s)$ denotes the total generation time spent by the agent to produce $s^\star$.
